\documentclass[12pt]{article}
\usepackage{fontspec}
\usepackage{xeCJK}
\usepackage{amsmath,amssymb}
\usepackage{tikz}
\usepackage{geometry}

\setCJKmainfont{Songti SC}
\setCJKsansfont{Heiti SC}
\setCJKmonofont{Songti SC}

\usetikzlibrary{automata,positioning,arrows}
\geometry{a4paper,margin=1in}
\renewcommand{\arraystretch}{1.2}

\begin{document}

\begin{center}
    {\LARGE\bfseries 编译原理作业 2\par}
    \vspace{1.2cm}
    {\large
    \begin{tabular}{rl}
    姓名：& 梁力航 \\
    学号：& 23336128 \\
    日期：& 2026年4月16日
    \end{tabular}\par}
\end{center}

\vspace{1cm}

\section*{第 1 题}
已知文法
\[
S \rightarrow aSb \mid ab
\]

\subsection*{1. 拓广文法与 LR(0) 项目}
拓广文法为
\[
S' \rightarrow S
\]
\[
S \rightarrow aSb \mid ab
\]

全部 LR(0) 项目为：
\begin{itemize}
    \item $S' \rightarrow \cdot S$
    \item $S' \rightarrow S \cdot$
    \item $S \rightarrow \cdot aSb$
    \item $S \rightarrow a \cdot Sb$
    \item $S \rightarrow aS \cdot b$
    \item $S \rightarrow aSb \cdot$
    \item $S \rightarrow \cdot ab$
    \item $S \rightarrow a \cdot b$
    \item $S \rightarrow ab \cdot$
\end{itemize}

\subsection*{2. LR(0) 项目集规范族与 DFA}
记产生式编号为
\begin{enumerate}
    \item $S \rightarrow aSb$
    \item $S \rightarrow ab$
\end{enumerate}

项目集规范族为：
\[
I_0 = \{S' \rightarrow \cdot S,\; S \rightarrow \cdot aSb,\; S \rightarrow \cdot ab\}
\]
\[
I_1 = goto(I_0,S) = \{S' \rightarrow S \cdot\}
\]
\[
I_2 = goto(I_0,a) = \{S \rightarrow a \cdot Sb,\; S \rightarrow a \cdot b,\; S \rightarrow \cdot aSb,\; S \rightarrow \cdot ab\}
\]
\[
I_3 = goto(I_2,S) = \{S \rightarrow aS \cdot b\}
\]
\[
I_4 = goto(I_2,b) = \{S \rightarrow ab \cdot\}
\]
\[
I_5 = goto(I_3,b) = \{S \rightarrow aSb \cdot\}
\]

且
\[
goto(I_2,a)=I_2
\]
所以 $I_2$ 上有一条关于终结符 $a$ 的自环。

对应 DFA 如下：
\begin{figure}[h]
    \centering
    \begin{tikzpicture}[shorten >=1pt,node distance=3cm,on grid,auto]
        \node[state,initial] (q0) {$I_0$};
        \node[state] (q1) [right=of q0] {$I_1$};
        \node[state] (q2) [below=of q0] {$I_2$};
        \node[state] (q3) [right=of q2] {$I_3$};
        \node[state] (q4) [below=of q3] {$I_4$};
        \node[state] (q5) [right=of q3] {$I_5$};

        \path[->]
        (q0) edge node {S} (q1)
             edge node[left] {a} (q2)
        (q2) edge[loop left] node {a} (q2)
             edge node {S} (q3)
             edge node[right] {b} (q4)
        (q3) edge node {b} (q5);
    \end{tikzpicture}
    \caption{文法 $G'$ 的 LR(0) 项目集 DFA}
\end{figure}

\subsection*{3. FOLLOW 集、SLR(1) 判断与分析表}
有
\[
FOLLOW(S')=\{\$\},\qquad FOLLOW(S)=\{b,\$\}
\]

归约项目只出现在 $I_4$ 与 $I_5$：
\begin{itemize}
    \item $I_4=\{S \rightarrow ab \cdot\}$，在 $FOLLOW(S)$ 上填入 $r2$；
    \item $I_5=\{S \rightarrow aSb \cdot\}$，在 $FOLLOW(S)$ 上填入 $r1$。
\end{itemize}

因此 SLR 分析表为：
\begin{center}
\begin{tabular}{|c|c|c|c|c|}
\hline
状态 & ACTION($a$) & ACTION($b$) & ACTION($\$$) & GOTO($S$) \\
\hline
0 & $s2$ &  &  & 1 \\
1 &  &  & acc &  \\
2 & $s2$ & $s4$ &  & 3 \\
3 &  & $s5$ &  &  \\
4 &  & $r2$ & $r2$ &  \\
5 &  & $r1$ & $r1$ &  \\
\hline
\end{tabular}
\end{center}

表中无冲突，所以该文法是 \textbf{SLR(1) 文法}。

\section*{第 2 题}
已知文法
\[
S \rightarrow E
\]
\[
E \rightarrow E+E \mid -E \mid id
\]

\subsection*{1. FIRST 集与 FOLLOW 集}
因为 $E$ 最终只能推出以 $-$ 或 $id$ 开始的串，所以
\[
FIRST(E)=\{-,id\},\qquad FIRST(S)=\{-,id\}
\]

又因为：
\begin{itemize}
    \item 由 $S \rightarrow E$，得 $\$ \in FOLLOW(E)$；
    \item 由 $E \rightarrow E+E$，得 $+ \in FOLLOW(E)$。
\end{itemize}
故
\[
FOLLOW(S)=\{\$\},\qquad FOLLOW(E)=\{+,\$\}
\]

\subsection*{2. 该文法为什么不是 SLR(1) 文法}
取 LR(0) 项目集中的关键状态为：
\[
I_0=\{S' \rightarrow \cdot S,\; S \rightarrow \cdot E,\; E \rightarrow \cdot E+E,\; E \rightarrow \cdot -E,\; E \rightarrow \cdot id\}
\]
\[
I_1=\{S \rightarrow E \cdot,\; E \rightarrow E \cdot +E\}
\]
\[
I_2=\{S' \rightarrow S \cdot\}
\]
\[
I_3=\{E \rightarrow - \cdot E,\; E \rightarrow \cdot E+E,\; E \rightarrow \cdot -E,\; E \rightarrow \cdot id\}
\]
\[
I_4=\{E \rightarrow id \cdot\}
\]
\[
I_5=\{E \rightarrow E+ \cdot E,\; E \rightarrow \cdot E+E,\; E \rightarrow \cdot -E,\; E \rightarrow \cdot id\}
\]
\[
I_6=\{E \rightarrow -E \cdot,\; E \rightarrow E \cdot +E\}
\]
\[
I_7=\{E \rightarrow E+E \cdot,\; E \rightarrow E \cdot +E\}
\]

在状态 $I_6$ 中，项目 $E \rightarrow -E \cdot$ 需要在 $FOLLOW(E)=\{+,\$\}$ 上归约，而项目 $E \rightarrow E \cdot +E$ 又要求在输入 $+$ 时移进，因此在符号 $+$ 上出现移进--归约冲突。

同理，在状态 $I_7$ 中，项目 $E \rightarrow E+E \cdot$ 与项目 $E \rightarrow E \cdot +E$ 也会在符号 $+$ 上产生移进--归约冲突。

所以该文法 \textbf{不是 SLR(1) 文法}。

\subsection*{3. 按给定冲突处理策略分析串 \texttt{-id+id\$}}
产生式编号记为：
\begin{enumerate}
    \item $S \rightarrow E$
    \item $E \rightarrow E+E$
    \item $E \rightarrow -E$
    \item $E \rightarrow id$
\end{enumerate}

题目给定策略为：
\begin{itemize}
    \item 移进--归约冲突时优先移进；
    \item 归约--归约冲突时优先采用编号较小的产生式归约。
\end{itemize}

本题实际只涉及移进--归约冲突，因此在状态 $I_6$ 遇到 $+$ 时取移进。

分析过程如下：
\begin{center}
\begin{tabular}{|c|c|c|c|c|c|}
\hline
步骤 & 状态栈 & 符号栈 & 输入串 & 动作 & 使用产生式编号 \\
\hline
1  & 0         & 空      & $-id+id\$$ & Shift $-$                   & - \\
2  & 0,3       & $-$     & $id+id\$$  & Shift $id$                  & - \\
3  & 0,3,4     & $-id$   & $+id\$$    & Reduce $E\to id$           & 4 \\
4  & 0,3,6     & $-E$    & $+id\$$    & Goto$(3,E)=6$              & - \\
5  & 0,3,6     & $-E$    & $+id\$$    & Shift $+$                   & - \\
6  & 0,3,6,5   & $-E+$   & $id\$$     & Shift $id$                  & - \\
7  & 0,3,6,5,4 & $-E+id$ & $\$$       & Reduce $E\to id$           & 4 \\
8  & 0,3,6,5,7 & $-E+E$  & $\$$       & Goto$(5,E)=7$              & - \\
9  & 0,3,6,5,7 & $-E+E$  & $\$$       & Reduce $E\to E+E$          & 2 \\
10 & 0,3,6     & $-E$    & $\$$       & Goto$(3,E)=6$              & - \\
11 & 0,3,6     & $-E$    & $\$$       & Reduce $E\to -E$           & 3 \\
12 & 0,1       & $E$     & $\$$       & Goto$(0,E)=1$              & - \\
13 & 0,1       & $E$     & $\$$       & Reduce $S\to E$            & 1 \\
14 & 0,2       & $S$     & $\$$       & Goto$(0,S)=2$              & - \\
15 & 0,2       & $S$     & $\$$       & Accept                      & - \\
\hline
\end{tabular}
\end{center}

因此，按题目给定的冲突处理规则，串 \texttt{-id+id\$} 可以被成功接受。

\end{document}
